% !TeX root = ../paper.tex
\section{Conclusion and Outlook}
\label{sec:conclusion}

This work addresses key challenges in clinical decision support for tropical and infectious diseases, where existing systems remain limited by insufficient multimodal integration, poor generalization to tropical infectious diseases, and lack of interpretability.

% To overcome these limitations, we implemented a multimodal diagnostic system leveraging large multimodal models, and systematically compared Multimodal \ac{RAG} (LightRAG + RAGAnything), domain-specific fine-tuning using NVIDIA NeMo, and a hybrid approach. We also established a comprehensive evaluation framework to assess diagnostic accuracy, reasoning quality, and reliability across diverse disease categories.

Our contributions to overcome these limitations include implementation and systematic comparative evaluation of three multimodal clinical decision-support approaches: Multimodal \ac{RAG} (LightRAG + RAGAnything), 
domain-specific fine-tuning using NVIDIA NeMo, and a hybrid architecture integrating both strategies. We further construct a standardized multimodal benchmark dataset and develop a 
unified fine-tuning dataset with structured diagnostic answers and step-by-step reasoning aligned with a USMLE-inspired MedCaseReasoning workflow. In addition, we design a rigorous 
evaluation framework leveraging RAGAS to assess diagnostic accuracy and reasoning quality, and implement an interactive LangGraph-based interface with step-level reasoning monitoring 
to enhance transparency and interpretability. 

When interpreting the results, we can observe that while the implemented models are not yet reliable as fully autonomous diagnostic clinical decision-support systems, 
their strong reasoning alignment and significant performance improvements through RAG and fine-tuning integration highlight their meaningful clinical potential. Specifically, the combination of external knowledge 
retrieval and enhanced reasoning capability positions these models as promising clinical decision-support tools, especially for generating structured and clinically relevant differential diagnoses. 
Future works should focus on further improving context precision and overall diagnostic accuracy by integrating an effective reranking module into the retrieval pipeline.
We additionally recommend expanding and refining the benchmark dataset to increase the disease diversity and clinical realism. Subsequently, exploring and adopting more clinically
meaningful evaluation metrics, such as top-k differential diagnosis accuracy, will better reflect real-world diagnostic performance.